% Resume - Vivekananda Swamy Mattam
% Technical founder / professor version

\documentclass[10pt,letterpaper]{article}
\usepackage[utf8]{inputenc}
\usepackage[margin=0.4in, top=0.25in, bottom=0.25in]{geometry}
\usepackage{enumitem}
\usepackage{hyperref}
\usepackage{titlesec}
\usepackage{xcolor}

% Link colors - blue
\hypersetup{
    colorlinks=true,
    linkcolor=blue,
    urlcolor=blue
}

% Formatting
\pagestyle{empty}
\setlength{\parindent}{0pt}
\setlength{\parskip}{0pt}

% Section formatting - bold uppercase with line
\titleformat{\section}{\normalsize\bfseries\uppercase}{}{0em}{}[\vspace{-4pt}\rule{\textwidth}{0.4pt}]
\titlespacing{\section}{0pt}{5pt}{2pt}

% Tight lists
\setlist[itemize]{leftmargin=0.15in, itemsep=0pt, parsep=0pt, topsep=1pt}

\begin{document}

%----------HEADING----------
\begin{center}
    {\Large \textbf{VIVEKANANDA SWAMY MATTAM}} \\
    \vspace{4pt}
    \small
    Robotics engineer focused on robust mobile autonomy, perception under sensor failure, and edge deployment for real robots. \\
    \vspace{3pt}
    +1 (949) 400-7126 \quad $\bullet$ \quad New York, NY \\
    \href{mailto:vm2677@nyu.edu}{vm2677@nyu.edu} \quad $\bullet$ \quad
    \href{https://linkedin.com/in/vivek-mattam-a8590b23a}{linkedin} \quad $\bullet$ \quad
    \href{https://github.com/vivekmattam02}{github} \quad $\bullet$ \quad
    \href{https://vivekmattam02.github.io}{portfolio}
\end{center}
\vspace{-4pt}

%----------EDUCATION----------
\section{Education}

\textbf{Master of Science in Mechatronics and Robotics Engineering} \hfill Sep 2024 -- May 2026 \\
\textit{New York University Tandon School of Engineering} \\
\textbf{Bachelor of Technology in Mechanical Engineering} \hfill Sep 2019 -- May 2023 \\
\textit{VNR VJIET, Hyderabad, India}
\vspace{-2pt}

%----------PROJECTS----------
\section{Selected Projects}

\href{https://github.com/vivekmattam02/ros2-ackermann-racing-navigation}{\textbf{Vortex: Robust Navigation with Edge-Deployed Perception}} $|$ \textit{Bell Labs Funded MS Project} \hfill Jan 2025 -- Present \\
\textit{Advisor: \href{https://wp.nyu.edu/aarab/}{Dr. Aliasghar Arab}}
\begin{itemize}
    \item Built a 1/10-scale autonomous Ackermann robot with Jetson Orin Nano, LiDAR, RGB-ToF sensing, VESC control, Teensy-based odometry, and custom ROS 2 nodes for motor control, TensorRT inference, and class-aware navigation.
    \item Distilled Depth Anything V3 and SAM2 into a 5.3M-parameter edge model for joint depth and segmentation, recovering 55\% more obstacle evidence when ToF depth degraded on reflective surfaces.
    \item Fused learned depth with LiDAR SLAM to cut median localization error by 73\% (4.63m to 1.23m APE) and ran the perception stack at 21.9 FPS on Jetson Orin Nano with TensorRT FP16.
\end{itemize}
\vspace{1pt}

\href{https://github.com/vivekmattam02/citywalker-earthrover}{\textbf{CityWalker-EarthRover Deployment}} $|$ \textit{VIP Self-Drive, AI4CE Lab} \hfill Sep 2024 -- Present \\
\textit{Faculty: \href{https://scholar.google.com/citations?user=YeG8ZM0AAAAJ}{Prof. Chen Feng}}
\begin{itemize}
    \item Deployed \href{https://ai4ce.github.io/CityWalker/}{CityWalker} zero-shot on FrodoBots EarthRover for sidewalk navigation, building the camera-to-control pipeline for preprocessing, waypoint prediction, coordinate transforms, and differential-drive tracking.
    \item Added depth barrier regularization with Depth-Anything-V2 at training time, reducing depth violations by 15\% and increasing minimum depth margin by 23\% while keeping navigation accuracy within 3\% of baseline.
    \item Led deployment and evaluation work within NYU's VIP Self-Drive team, including Earth Rover Challenge 3 preparation and field testing.
\end{itemize}
\vspace{1pt}

\href{https://github.com/vivekmattam02/spot}{\textbf{Autonomous Person Following on Boston Dynamics Spot}} \hfill Jan 2025 -- Present
\begin{itemize}
    \item Built a real-time visual-servoing system that uses YOLOv8 detections to estimate lateral, distance, and pitch errors and generate smooth 10 Hz velocity commands for person following.
    \item Integrated Spot SDK, ZED SDK, OpenCV, CUDA, and Docker into an on-robot perception-and-control stack, including body-pitch compensation for tracking on stairs and slopes.
\end{itemize}
\vspace{1pt}

\href{https://github.com/vivekmattam02/LAC_26}{\textbf{Lunar Autonomy Challenge -- Perception-Aware Navigation}} \hfill Summer 2025
\begin{itemize}
    \item Built an autonomous stack for a simulated lunar rover without GPS or prior maps, combining ORB-SLAM3 pose estimation, stereo depth, voxel and height mapping, A* planning, and pure-pursuit plus PID control.
    \item Propagated stereo uncertainty through mapping and planning so the rover actively avoided shadows, crater rims, and texture-poor terrain instead of treating depth as ground truth.
\end{itemize}
\vspace{1pt}

\href{https://github.com/Nishant-ZFYII/rob6323_go2_project}{\textbf{Reinforcement Learning for Quadruped Locomotion}} $|$ \textit{NVIDIA Isaac Lab} \hfill Sep 2024 -- Dec 2024
\begin{itemize}
    \item Trained a Unitree Go2 locomotion policy with PPO using reward shaping, Raibert heuristic gait coordination, and actuator-friction domain randomization for improved sim-to-real robustness.
    \item Achieved nearly 2x baseline velocity-tracking performance on flat and rough terrain.
\end{itemize}

%----------EXPERIENCE----------
\section{Experience}

\href{https://wp.nyu.edu/aarab/}{\textbf{Course Assistant -- Autonomous Mobile Robots}} $|$ \textit{NYU Tandon, Prof. Aliasghar Arab} \hfill Sep 2024 -- Present
\begin{itemize}
    \item Developed lecture notes, assignments, diagrams, and the \href{https://vivekmattam02.github.io/amr-notes/}{course website} for a graduate robotics course; the material now serves as reusable infrastructure for instruction in controls and motion planning.
\end{itemize}
\vspace{1pt}

\textbf{Robotics Engineering Intern} $|$ \textit{Xmachines -- Agricultural Robotics Startup} \hfill Feb 2024 -- Aug 2024
\begin{itemize}
    \item Built a Flask-based telemetry interface that streamed IMU and camera data for debugging robot state during field tests.
    \item Developed a computer-vision sorting tunnel for a weeder robot and supported ROS-based motion planning plus embedded integration on Jetson Orin Nano, Arduino, and ESP32.
\end{itemize}

%----------SKILLS----------
\section{Skills}

\begin{itemize}[leftmargin=0in, label={}, itemsep=0pt]
    \item \textbf{Languages:} Python, C++, C, MATLAB
    \item \textbf{Robotics:} ROS, ROS 2, Nav2, SLAM Toolbox, ORB-SLAM3, Visual Odometry, EKF, MPC, MPPI, A*, Pure Pursuit
    \item \textbf{Perception \& ML:} PyTorch, TensorRT, ONNX, OpenCV, YOLOv8/v11, ByteTrack, CosPlace, SuperGlue, RANSAC, ICP, Open3D
    \item \textbf{Systems:} Jetson Orin Nano, Boston Dynamics Spot, TurtleBot3, Gazebo, Isaac Lab, Docker, Git, Linux, Arduino, ESP32
\end{itemize}

\end{document}
